\begin{proposition}[The special supercharges and the supercharges]
\label{prop:c11-S}
Let $\eta\in\Sigma^*$ and $b\in V$, so that $\rho(b)\eta\in\Sigma$.
\begin{itemize}
\item (i) Graph by graph, the position part along $b$ of the two-body kernel of $(X_a,X_b)$ along $\eta$ is the two-body kernel of $(X_a,X_b)$ along $\rho(b)\eta$; its translation-invariant part is computed with $S^{(1)}(\eta)$ at the spinor $\rho(x-x_O)\eta$, $x_O$ the point of the density, together with the term $-\int(\rho(A|_{W_{\mathbb R}^\perp})\eta,\lambda^*)$.
\item (ii) Under the hypotheses of Proposition \ref{prop:c7-select} and Corollary \ref{cor:c8-homotopy}, on cohomology, the position part along $b$ of $S_1$ along $\eta$ is $Q_1$ along $\rho(b)\eta$ minus the commutator of $\{S^{(1)}(\eta),-\}$ with $P_1$.
\item (iii) On a family on which $P_1\cong0$, in particular, under the hypotheses of Theorem \ref{thm:c11-arity2}, on the families of Definition \ref{def:c11-families} in the translation-covariant scheme at $\theta=\theta_\lambda=\theta_V=-\tfrac13$, the position part of $S_1$ along $\eta$ is $Q_1$ along $\rho(b)\eta$ on cohomology, and $S_1$ along $\eta$ is fixed by $Q_1$ along $\Sigma$ up to the term of Proposition \ref{prop:c7-select} linear in ${\not d}(\rho(x)\eta)=4\eta$. On the scalar bilinears this term is absent (Corollary \ref{cor:c7-S}), and $S_1$ along $\eta$ is, on cohomology, the combination of the functionals of Corollary \ref{cor:c7-S} with the coefficients of $Q_1$ along $\Sigma$ in Corollary \ref{cor:c7-Q}, given by Theorem \ref{thm:c11-bilinears} and Proposition \ref{prop:c11-SKbil} (i).
\item (iv) In a reflection-compatible scheme, on the family formed by a traceless scalar bilinear $O$ at the origin and its descendants $\{S^{(1)}(\xi),O\}$, on which $Q_1=c_Q\{S^{(1)}(\xi),-\}$ along the unit generator $Q^i_\alpha$ (Theorem \ref{thm:c11-bilinears} (ii)) and $S_1=c_S\{S^{(1)}(\eta),-\}$ along $S^\alpha{}_i$, Corollary \ref{cor:c8-gram} reads
\[
c_S=\overline{c_Q}+g_{5/2}-g_2,
\]
with $g_w$ the eigenvalue of $\mathrm{Gram}_{\rm tree}^{-1}\mathrm{Gram}_{\rm 1}$ in Weyl weight $w$ on the family; the outer direction changes $c_Q$ by $-a$ and $c_S$ by $+a$ and preserves the relation for $\bar a=-a$.
\end{itemize}
\end{proposition}

\begin{proof}
(i) By \eqref{eq:c11-Seta}, $S^{(1)}(\eta)=\int x^\mu s_\mu+\int s_0$, where $\int s_\mu$ is $S^{(1)}$ along the constant spinor $\rho(e_\mu)\eta\in\Sigma$. Insert $x^\mu=x_O^\mu+(x-x_O)^\mu$ and argue as in the proof of Proposition \ref{prop:c11-position} (i): on plane waves the factor $x^\mu$ is $-i\partial/\partial r^\mu$ of the momentum injected at the vertex, and its action on the momentum delta function is $-i$ times the graph with $\int s_\mu$; the same holds for the special edges, the vertices $S^{(1)}(\eta)$ on the cycle and the counterterm vertices, which are the defect insertions of Proposition \ref{prop:c10-arity1} (Theorem \ref{thm:c11-oddsimplex} (ii)). (ii) Corollary \ref{cor:c8-homotopy} (ii) at $(\eta,P_b)$: by Proposition \ref{prop:c2-special}, $[P_b,\eta]=\rho(b)\eta$, so $[\eta,P_b]=-\rho(b)\eta$, and $S^{(2)}(\eta,P_b)=0$ (Definition \ref{def:c2-S2}). At order $\hbar$ the identity states that the commutator of $S_1$ along $\eta$ with $\{S^{(1)}(P_b),-\}$, plus the commutator of $\{S^{(1)}(\eta),-\}$ with $P_1$, minus $Q_1$ along $\rho(b)\eta$, is a graded commutator with $\Lambda^{(0)}$, which induces $0$ on cohomology; at order $\hbar^0$ the same identity is the statement of (i) for the classical transformation, which fixes the signs. By Proposition \ref{prop:c10-bound1}, $S_1$ along $\eta$ is at most linear in the position, and Lemma \ref{lem:c11-canonical} identifies the first commutator with its position part. (iii) By Proposition \ref{prop:c7-select} the representative of $S_1$ along $\eta$ at a point $x_0$ is a term linear in the value $\rho(x_0)\eta$ plus a term linear in $4\eta$; the first has the position part given by its value at $\rho(b)\eta$, and as $\eta$ and $b$ vary, $\rho(b)\eta$ spans $\Sigma$, so by (ii) it is the representative of $Q_1$ evaluated at $\rho(x_0)\eta$ on cohomology. (iv) Write $X_0+\hbar X_1$ for the matrix of $-\Lambda^{(1)}(Q^i_\alpha)$ and $Y_0+\hbar Y_1$ for that of $-\Lambda^{(1)}(-S^\alpha{}_i)$ (Theorem \ref{thm:c8-dag} (i)): $\hbar X_1=c_QX_0$ and $\hbar Y_1=c_SY_0$, and $Y_0$ maps Weyl weight $\tfrac52$ to Weyl weight $2$. Corollary \ref{cor:c8-gram} gives $\hbar Y_1=\overline{c_Q}\,X_0^{\rm adj}+\hbar[Y_0,\mathrm{Gram}_{\rm tree}^{-1}\mathrm{Gram}_{\rm 1}]=\big(\overline{c_Q}+g_{5/2}-g_2\big)Y_0$ on the descendants, with $Y_0=X_0^{\rm adj}$. The outer direction: Proposition \ref{prop:c11-freedom} (ii) and Corollary \ref{cor:c8-gram}.
\end{proof}

\begin{corollary}[$K_1$ in the translation-covariant scheme]
\label{cor:c11-K}
Under the hypotheses of Theorems \ref{thm:c8-adjoint} and \ref{thm:c11-arity2}, on a family of Definition \ref{def:c11-families} at the origin of one length, in the translation-covariant scheme at $\theta=\theta_\lambda=\theta_V=-\tfrac13$, $K_1$ along $K_{\bar b}$ equals $\hbar\big[\{S^{(1)}(K_{\bar b}),-\},\mathrm{Gram}_{\rm tree}^{-1}\mathrm{Gram}_{\rm 1}\big]$, the statement $K=P^\dagger$ of Corollary \ref{cor:c8-gram}. For equal coefficients $\theta_P$ along $P_b$, $K_1$ along $K_{\bar b}$ is $(\theta-\theta_P)$ units times $\{S^{(1)}(K_{\bar b}),-\}$ plus the same commutator, with the radial Gram matrix of Definition \ref{def:c8-inner} in that scheme. On a scalar bilinear at the origin both sides vanish: the defect cycles give $K_1(O)=0$ there (Proposition \ref{prop:c11-SKbil} (ii)), and the commutator vanishes on a primary at the origin on which $\mathrm{Gram}_{\rm tree}^{-1}\mathrm{Gram}_{\rm 1}$ is diagonal. On its descendants $\partial_\mu O$ the relation equates the intrinsic kernels $\Upsilon^0_{K_b}$ of the descendant pairs with the difference of the two-point normalisations in Weyl weights $w_O+1$ and $w_O$ times the classical transformation.
\end{corollary}

\begin{proof}
Corollary \ref{cor:c11-D1fam} gives $P_1\cong0$, resp. $P_1\cong(\theta-\theta_P)$ units times $\{S^{(1)}(P_b),-\}$ with a real coefficient; Corollary \ref{cor:c8-gram} gives the two relations. At the origin $v(0)=0$ and $(b,0)=0$, so the formula of Proposition \ref{prop:c11-SKbil} (ii) vanishes; $\{S^{(1)}(K_{\bar b}),-\}$ lowers the Weyl weight by one and annihilates a primary at the origin.
\end{proof}

\begin{proposition}[The relations at order $\hbar$]
\label{prop:c11-relations}
Write $\Lambda^{(1)}(x)=\Lambda^{(1)}_0(x)+\hbar\Lambda^{(1)}_1(x)$ with $\Lambda^{(1)}_0(x)=\pm\{S^{(1)}(x),-\}$ (Theorem \ref{thm:c8-L} (iv)), so that $\hbar\Lambda^{(1)}_1(x)O$ is the component of $O$ along $x$, and $\Lambda^{(0)}=s+\hbar L_1^{(0)}$. Let $\mathfrak F$ be a family of colour-singlet local functionals of degree $0$, stable under $\Lambda$ and under the classical transformations, with $L_1^{(0)}=0$ on $\mathfrak F$ for a finite part $\psi$ as in Definition \ref{def:c10-hyp} (ii), and let $O\in\mathfrak F$ with $sO=0$.
\begin{itemize}
\item (R1) $s\big(\hbar\Lambda^{(1)}_1(x)O\big)=0$ for every generator $x$; along $\xi$, $s\,Q_1(O)=0$, an identity of local functionals.
\item (R2) For $\xi,\xi'\in\Sigma$ with $b=\Gamma(\xi,\xi')|_V$, $[\Lambda^{(1)}_1(\xi),\Lambda^{(1)}_0(\xi')]O+[\Lambda^{(1)}_0(\xi),\Lambda^{(1)}_1(\xi')]O+\Lambda^{(1)}_1(P_b)O=\Lambda^{(0)}(\chi)$, with $\chi$ local of degree $-1$, built from the arity-two component of $O$ along $(\xi,\xi')$; for $\xi,\xi'\in S_+\otimes W$, $b=0$.
\item (R3) For $\xi\in\Sigma$ and $\eta\in\Sigma^*$ with $[\xi,\eta]=\tfrac12(\xi,\eta)\Delta+M_{\xi\eta}+R_{\xi\eta}$ (Proposition \ref{prop:c2-special}), $M_{\xi\eta}\in\mathfrak{so}(V)$, $R_{\xi\eta}\in\mathfrak{so}(W_{\mathbb R}^\perp)$,
\[
[\Lambda^{(1)}_1(\xi),\Lambda^{(1)}_0(\eta)]O+[\Lambda^{(1)}_0(\xi),\Lambda^{(1)}_1(\eta)]O+\tfrac12(\xi,\eta)\,D_1(O)+M_1(O)+R_1(O)=\Lambda^{(0)}(\chi'),
\]
the components $M_1$, $R_1$ taken along $M_{\xi\eta}$, $R_{\xi\eta}$.
\item (R4) For a rank-one $\xi$, $Q^\xi_0Q^\xi_1+Q^\xi_1Q^\xi_0=0$ on every local functional (Theorem \ref{thm:c8-twist}).
\item (R5) With $[K_b,P_{b'}]=(b,b')\Delta+M_{bb'}$ (Remark \ref{rem:c7-KDelta}), $[\Lambda^{(1)}_1(K_b),\Lambda^{(1)}_0(P_{b'})]+[\Lambda^{(1)}_0(K_b),\Lambda^{(1)}_1(P_{b'})]+(b,b')\hbar\Lambda^{(1)}_1(\Delta)+\hbar\Lambda^{(1)}_1(M_{bb'})\cong0$ on cohomology; with Theorem \ref{thm:c11-oddsimplex} (iii) it states that the parts of $K_1$ multiplying the position are those of \eqref{eq:c11-Kdec}.
\item (R6) For $\eta_1,\eta_2\in\Sigma^*$ with $[\eta_1,\eta_2]=-K_b$, $b=\Gamma(\eta_1,\eta_2)|_V$, at a finite $T$, $K_1$ along $K_b$ is, on cohomology, $[\Lambda^{(1)}_1(\eta_1),\Lambda^{(1)}_0(\eta_2)]+[\Lambda^{(1)}_0(\eta_1),\Lambda^{(1)}_1(\eta_2)]$.
\end{itemize}
In (R2), (R3), (R5) and (R6) the right sides vanish on cohomology. For $O$ without letters $\lambda$ and $c$, such as the scalar bilinears and the families of scalars, the parts of order $\hbar^0$ of $\chi$, $\chi'$ vanish, the right sides are $s$-exact, and their parts with two field letters and no antifield vanish on free on-shell plane waves, $|p|^2=0$ on the scalars and the gauge field, $\rho(p)\lambda=0$ and $(p,a)=0$ on the polarisations $a$ of $A|_V$. On the scalar bilinears at $\theta=-\tfrac13$, (R3) contains $D_1(O)=0$ on the traceless part and $-48$ units times $O$ on the trace (Corollary \ref{cor:c9-ident}), the Konishi relation between $Q_1$, $S_1$ and the anomalous dimension, and (R2) with $\xi\in S_+\otimes W$, $\xi'\in S_-\otimes W^*$ contains $P_1(O)=-(\tfrac13+\theta_P)$ units times $\partial_bO$ (Theorem \ref{thm:c9-P1}).
\end{proposition}

\begin{proof}
Corollary \ref{cor:c8-homotopy} (i) at order $\hbar$ reads $[s,\Lambda^{(1)}_1(x)]+[L_1^{(0)},\Lambda^{(1)}_0(x)]=0$; on $O$ with $sO=0$, $L_1^{(0)}O=0$ and $L_1^{(0)}\{S^{(1)}(x),O\}=0$ (stability of $\mathfrak F$) this is (R1). Corollary \ref{cor:c8-homotopy} (ii) at order $\hbar$ for the pairs $(\xi,\xi')$, $(\xi,\eta)$, $(K_b,P_{b'})$, $(\eta_1,\eta_2)$, with the brackets of Proposition \ref{prop:c2-special} and Remark \ref{rem:c7-KDelta}, has the right side $[\Lambda^{(0)},H]$ with $H$ built from $\Lambda^{(2)}$, which on $O$ with $\Lambda^{(0)}O=0$ is $\Lambda^{(0)}(HO)$; the part of order $\hbar^0$ of $H$ is $\pm\{S^{(2)}(x_1,x_2),-\}$, and $S^{(2)}$ has the legs $(\lambda^*,\lambda^*)$ and $(A,c^*)$ (Definition \ref{def:c2-S2}), whose bracket with a functional of the fields $A$ alone vanishes. (R4) is Theorem \ref{thm:c8-twist}. For the vanishing on shell: the part of $s\chi$ with two field letters is $Q\chi_2+\{I_3,\chi_1\}$, with $\chi_k$ the part of $\chi$ with $k$ letters; $\chi_1$ is a single adjoint letter contracted with a constant tensor, zero for a colour singlet; $\chi_2$ of degree $-1$ has one antifield and one field, or $c^*$ and $c$, and $Q$ substitutes $A^*\mapsto-\tfrac12d^*dA$, $\lambda^*\mapsto{\not d}\lambda$, $c^*\mapsto d^*(A^*|_V)$, $A\mapsto-dc$ (Proposition \ref{prop:c1-Q}), so the part of $Q\chi_2$ without antifields and ghosts is a free field equation times a field. The values on the bilinears are Corollary \ref{cor:c9-ident} and Theorem \ref{thm:c9-P1}, with $R_1=0$ on them (Remark \ref{rem:c7-internal}).
\end{proof}

